\subsection{Mathematical Foundations and Structural Syntax}

A \texttt{set} is a hash-backed uniqueness domain: insertion collapses equal elements, iteration exposes current members without positional semantics, and only hashable objects may enter the table.

\subsubsection{Set Architectures: The Underlying Dictionary Engine}

CPython implements \texttt{set} as a hollow variant of \texttt{PyDictObject}. The same open-addressed hash table, probe loop, and entry layout used for dictionaries stores set elements as \emph{keys only}; every value slot points at a single shared dummy object (\texttt{PySet\_Dummy} in C sources) rather than a distinct payload.

Consequences follow directly from that reuse:

\begin{itemize}
    \item Elements inherit dict-key hashability rules: immutable, stable \texttt{\_\_hash\_\_} and \texttt{\_\_eq\_\_} for the object's lifetime in the table.
    \item Membership testing executes the identical probe sequence as dict key lookup---hash compare first, \texttt{\_\_eq\_\_} second---with amortized $O(1)$ cost under normal load.
    \item Resize, deletion-marker, and perturbation logic are shared engineering; a \texttt{set} is not a separate ad-hoc container but a value-stripped \texttt{dict}.
\end{itemize}

Surface syntax (\texttt{\{a, b\}}, \texttt{set()}, comprehensions) constructs the same \texttt{PySetObject} header regardless of literal form; \texttt{\{\}} remains an empty \texttt{dict} by historical grammar, not an empty set.
